%% kapa: exam=1 round=27 year=2016 period=2 form=A subject=appraisal_law subject_ko=감정평가관계법규 seq=5 number=5
%% source: Q-Net 감정평가사 기출문제 (cst00309) artlSeq=5169402 file=extracted/1008/5169402/제27회 감정평가사 1차 문제/제27회 감정평가사 2교시문제(A형).pdf
%% conversion: kordoc 4.24.51 -> tools/gen1.py (auto); review=proofread-sonnet; flags=none
\begin{question}{5}
국토의 계획 및 이용에 관한 법령상 도시ㆍ군기본계획을 수립하지 않을 수 있는 지방자치단체는? (단, 수도권은「수도권정비계획법」상의 수도권을 의미함)
\begin{choices}
  \item 수도권에 속하는 인구 10만명 이하인 군
  \item 수도권에서 광역시ㆍ특별시와 경계를 같이하는 인구 10만명 이하인 시
  \item 수도권외 지역에서 광역시와 경계를 같이하지 아니하는 인구 10만명 이하인 시
  \item 관할구역 일부에 대하여 광역도시계획이 수립되어 있는 시로서 광역도시계획에 도시ㆍ군기본계획의 내용이 모두 포함되어 있는 시
  \item 관할구역 전부에 대하여 광역도시계획이 수립되어 있는 군으로서 광역도시계획에 도시ㆍ군기본계획의 내용이 일부 포함되어 있는 군
\end{choices}
\end{question}
\answer{3}
